% Shared exercise chunk: l4-multistep
\begin{exercisechunk}
\item \textbf{Computing further-ahead predictions.}
\label[exercise]{ex:multistep}
Use the autoregressive model in \cref{eq:multistep-marginal}. Fix
\(x_{\le t}\), \(k\ge2\) with \(t+k\le T\), and \(a\in[N]\).
\begin{enumerate}[(a),beginpenalty=10000]
\item Derive \cref{eq:multistep-marginal}.
Count its terms and the number of conditional evaluations if each term
is evaluated separately. Explain why sharing prefixes still leaves
exponentially many distinct prefixes in a general model.
\item Which conditional probability is trained if the observed
intervening tokens are supplied as context instead of being marginalized?
\item Suppose \(t\ge1\) and the model is first-order Markov on \([N]\)
with known transition matrix \(K\). Give a recursion that computes the
entire \(k\)-step distribution in \(O(kN^2)\) arithmetic operations.
Why can this recursion merge histories?
\item A separate head \(h_{\theta,k}(\,\cdot\mid x_{\le t})\) can be
trained from pairs \((X_{\le t},X_{t+k})\). State the additional
equality needed for its probabilities to be the original sequence
model's \(k\)-step marginals. Does perfect prediction on the observed
prefix--target pairs imply this equality? Give an example or a proof.
\end{enumerate}
\end{exercisechunk}
